\documentclass[./../main.tex]{subfiles}

\begin{document}
\section{chapter4 动态规划}
当子问题不独立时，如果能够保存已解决的子问题的答案，而在需要时再找出已求得的答案，就可以避免大量重复计算，从而得到多项式时间算法。而动态规划算法只解决每个子问题一次，然后将其答案保存起来，从而避免了每次遇到子问题时重新计算答案的工作。即动态规划通过空间开销来节省由于重复计算带来的时间开销。

\textbf{最优性原理}：如果问题的最优解所包含的子问题的解也是子问题的最优解，就称该问题具有最优子结构，即满足最优性原理（优化原则）。

并不是所有优化问题都适用优化原则，有时对某个子问题的解不一定最优，但延伸成整个问题的解时反而最优。

在证明一个问题满足最优子结构性质时，可以考虑反证法，我们后面在0-1背包问题中具体展开。

\subsection{最长公共子序列}
要证明最长公共子序列问题满足最优子结构性质，设$X=\langle x_1,\ldots,x_m \rangle$和$Y=\langle y_1,\ldots,y_n \rangle$的LCS为$Z=\langle z_1,\ldots,z_k \rangle$。当$x_m = y_n$时必有$z_k = x_m$，否则可将$x_m$追加到$Z$得到更长子序列导致矛盾，此时$Z_{k-1}$必须是$X_{m-1}$和$Y_{n-1}$的LCS，若存在更长的$W$则$W \oplus x_m$将形成更长公共子序列；当$x_m \neq y_n$时，若$z_k \neq x_m$则$Z$是$X_{m-1}$和$Y$的LCS，若$z_k \neq y_n$则$Z$是$X$和$Y_{n-1}$的LCS，否则可通过选择更长子序列构造矛盾。由此LCS的最优解总由子问题最优解构成，满足最优子结构性质。

$s_{[-1]}$ 表示序列$s$的最后一个元素，$s_{[:-1]}$ 表示$s$去掉最后一个元素后的子序列
$$
\operatorname{LCS}\left(s_{1}, s_{2}\right)=
\begin{cases} 
0, & \text{if } s_{1} \text{ or } s_{2} \text{ is empty} \\
\operatorname{LCS}\left(s_{1_{[:-1]}}, s_{2_{[:-1]}}\right) + 1, & \text{if } s_{1_{[-1]}} = s_{2_{[-1]}} \\
\max\left(
    \operatorname{LCS}\left(s_{1_{[:-1]}}, s_{2}\right),
    \operatorname{LCS}\left(s_{1}, s_{2_{[:-1]}}\right)
\right), & \text{otherwise}
\end{cases}
$$

\begin{lstlisting}
while (cin >> s1 >> s2){
    memset(dp, 0, sizeof(dp));
    int n1 = strlen(s1), n2 = strlen(s2);
    for (int i = 1; i <= n1; ++i)
        for (int j = 1; j <= n2; ++j)
            if (s1[i - 1] == s2[j - 1])
                dp[i][j] = dp[i - 1][j - 1] + 1;
            else
                dp[i][j] = max(dp[i - 1][j], dp[i][j - 1]);
    cout << dp[n1][n2] << endl;
}
\end{lstlisting}
\subsection{最长上升子序列}
先说一种最普遍的方法，因为所求为子序列，所以这道题很容易想到一种线性动态规划。我们需要求最长上升子序列，为了上升我们肯定要知道我们当前阶段最后一个元素为多少，为了最长我们还要知道当前我们的序列有多长。我们可以用前者来充当第一维描述：设 $F[i]$ 表示以 $A[i]$ 为结尾的最长上升子序列的长度，为了保证元素单调递增我们肯定只能从 $i$ 前面的且末尾元素比 $A[i]$ 小的状态转移过来：

$$
F[i] =
\max_{0 \leq j < i, A[j] < A[i]} \{F[j] + 1\}
$$

初始值为 $F[0] = 0$，而答案可以是任何阶段中只要长度最长的那一个，所以我们边转移边统计答案。

复杂度：$O(n^{2})$

贪心加二分查找。

我们肯定要知道我们当前阶段最后一个元素为多少，还有当前我们的序列有多长。前两种方法都是用前者做状态，我们为什么不可以用后者做状态呢？

设 $F[i]$ 表示长度为$i$的最长上升子序列的末尾元素的最小值，我们发现这个数组的权值一定单调不降（仔细想一想，这就是我们贪心的来由）。于是我们按顺序枚举数组$A$，每一次对$F[]$数组二分查找，找到小于$A[i]$的最大的$F[j]$，并用它将$F[j+1]$更新。

期望复杂度：$O(n\log n)$

\begin{lstlisting}
int lengthOfLIS(vector<int>& nums) {
    vector<int> tails;  // tails[i]表示长度为i+1的LIS的最小末尾元素
    for (int num : nums) {
        // 使用二分查找找到第一个大于等于num的元素位置
        auto it = lower_bound(tails.begin(), tails.end(), num);
        if (it == tails.end()) {
            // 如果num比所有元素都大，可以扩展LIS长度
            tails.push_back(num);
        } else {
            // 否则替换第一个大于等于num的元素，保持tails数组的最小性
            *it = num;
        }
    }
    return tails.size();  // tails的大小就是LIS的长度
}
\end{lstlisting}

要同时求出最长公共子序列的长度和序列，首先初始化二维数组$c[m+1][n+1]$（存储长度）和$b[m+1][n+1]$（存储方向），将$c$的第一行和第一列初始化为0。然后遍历字符串$X$和$Y$：当$X[i-1]=Y[j-1]$时，令$c[i][j]=c[i-1][j-1]+1$并在$b[i][j]$标记"↖"；否则比较$c[i-1][j]$和$c[i][j-1]$，若$c[i-1][j]≥c[i][j-1]$则令$c[i][j]=c[i-1][j]$且标记"↑"，否则令$c[i][j]=c[i][j-1]$且标记"←"。最终$c[m][n]$即为LCS长度。回溯时从$b[m][n]$开始：若标记为"↖"则递归$b[i-1][j-1]$并记录$X[i-1]$；若为"↑"则递归$b[i-1][j]$；若为"←"则递归$b[i][j-1]$。最后将记录的字符反转即得LCS序列。

\begin{algorithm}[t]
    \KwIn{字符串$X = x_1x_2\cdots x_m$，字符串$Y = y_1y_2\cdots y_n$}
    \KwOut{LCS长度$c[m][n]$，方向矩阵$b[1..m][1..n]$}
    \BlankLine
    \For{$i \gets 0$ \KwTo $m$}{
        $c[i][0] \gets 0$\;
    }
    \For{$j \gets 0$ \KwTo $n$}{
        $c[0][j] \gets 0$\;
    }
    \For{$i \gets 1$ \KwTo $m$}{
        \For{$j \gets 1$ \KwTo $n$}{
            \eIf{$x_i == y_j$}{
                $c[i][j] \gets c[i-1][j-1] + 1$\;
                $b[i][j] \gets "\nwarrow"$\;
            }{
                \eIf{$c[i-1][j] \geq c[i][j-1]$}{
                    $c[i][j] \gets c[i-1][j]$\;
                    $b[i][j] \gets "\uparrow"$\;
                }{
                    $c[i][j] \gets c[i][j-1]$\;
                    $b[i][j] \gets "\leftarrow"$\;
                }
            }
        }
    }
    \Return $(c[m][n], b)$\;
    \caption{LCS-Length}
\end{algorithm}

\begin{algorithm}[t]
    \KwIn{方向矩阵$b[1..m][1..n]$，字符串$X$，位置$(i,j)$}
    \KwOut{LCS序列}
    \BlankLine
    \If{$i == 0$ \textbf{or} $j == 0$}{
        \Return\;
    }
    \uIf{$b[i][j] == "\nwarrow"$}{
        PRINT-LCS($b$, $X$, $i-1$, $j-1$)\;
        输出$x_i$\;
    }
    \uElseIf{$b[i][j] == "\uparrow"$}{
        PRINT-LCS($b$, $X$, $i-1$, $j$)\;
    }
    \Else{
        PRINT-LCS($b$, $X$, $i$, $j-1$)\;
    }
    \caption{PRINT-LCS}
\end{algorithm}

\subsection{矩阵链乘法}
给定$n$个矩阵的链$\langle A_1, A_2, \ldots, A_n \rangle$，其中矩阵$A_i$的维度为$p_{i-1} \times p_i$

设一个由$n$个矩阵组成的序列的可选的括号化方案为$P(n)$。则$P(n)$递归式(卡特兰数)为：

$$
P(n)=\begin{cases} 
1 & n=1\\
\sum\limits_{k=1}^{n-1} P(k)P(n-k) & n\geq 2
\end{cases}
$$

其闭合形式解为：
$$
P(n)=\frac{1}{n}\binom{2n-2}{n-1}=\frac{(2n-2)!}{n((n-1)!)^2}\approx\frac{4^n}{4\sqrt{\pi}n^{1.5}}
$$


我们的求解过程分以下几步：
\begin{enumerate}
    \item 描述最优解的结构特征。
    \item 递归定义最优解的值。
    \item 计算最优解的值。
    \item 根据计算信息构造最优解。
\end{enumerate}

定义$A_{i:j}$为$A_iA_{i+1}\cdots A_j$的乘积，其中$i \leq j$。如果$A_{i:j}$可以继续分割（即$i < j$），利用括号化分割成$A_{i:k}$和$A_{k+1:j}$，其中$i \leq k < j$。


$A_{i:j}$的括号化后的计算代价为：
$$
\text{计算代价}(A_{i:j}) = \text{计算代价}(A_{i:k}) + \text{计算代价}(A_{k+1:j}) + \text{相乘代价}(A_{i:k}, A_{k+1:j})
$$

关键性质：
\begin{itemize}
    \item 子问题的最优括号化方案包含在原问题的最优括号化方案中
    \item 通过求出所有子问题的最优解，就能组成原问题的最优解
    \item 在求解子问题的最优分割方案时，需要对所有可能的分割点$k$进行考察
\end{itemize}

假设矩阵链乘问题\textbf{不}满足最优子结构性质，即存在某个子问题的最优解不是构成原问题最优解的一部分。

考虑矩阵链$A_1A_2\cdots A_n$的最优括号化方案，设其在位置$k$处分割：
$$ A_{1:n} = (A_{1:k})(A_{k+1:n}) $$

根据假设，存在以下两种情形之一：
\begin{enumerate}
    \item $A_{1:k}$不是子问题$A_{1:k}$的最优括号化方案
    \item $A_{k+1:n}$不是子问题$A_{k+1:n}$的最优括号化方案
\end{enumerate}

以第一种情形为例（第二种情形对称），设存在更优的$A_{1:k}^*$使得：
$$ m^*[1,k] < m[1,k] $$

则可构造新的括号化方案：
$$ A_{1:n}' = (A_{1:k}^*)(A_{k+1:n}) $$

其总计算代价为：
\begin{align*}
m'[1,n] &= m^*[1,k] + m[k+1,n] + p_0p_kp_n \\
&< m[1,k] + m[k+1,n] + p_0p_kp_n \\
&= m[1,n]
\end{align*}

这与$m[1,n]$是最优解的初始假设矛盾。因此，原假设不成立，矩阵链乘必须满足最优子结构性质。

定义$m[i,j]$表示计算$A_{i:j}$的乘积所需标量乘法的最少次数。计算$A_{1}A_{2}\cdots A_{n}$的乘积所需标量乘法的最少次数可表示为$m[1,n]$。

我们可以递归求解$m[i,j]$：
\begin{itemize}
\item 当$i=j$时，$A_{i,i}=A_{i}$，$m[i,i]=0$，其中$i=1,2,\ldots,n$
\item 当$i<j$时，利用最优括号化方案：
$$m[i,j]=m[i,k]+m[k+1,j]+p_{i-1}p_{k}p_{j},\quad \text{其中}i\leq k<j$$
\end{itemize}

分段函数形式表示为：
$$
m[i,j]=\begin{cases}
0 & i=j\\
\min\limits_{i\leq k<j}\big\{m[i,k]+m[k+1,j]+p_{i-1}p_{k}p_{j}\big\} & i<j
\end{cases}
$$

定义$s[i,j]$表示$A_{i:j}$的最优分割方案中的值$k$，即满足：
$$
m[i,j]=m[i,k]+m[k+1,j]+p_{i-1}p_{k}p_{j}
$$

\begin{algorithm}[t]
    \KwIn{矩阵维度序列$p[0..n]$}
    \KwOut{最小乘法次数$m[n,n]$和最优分割方案$s[1..n-1,2..n]$}
    \BlankLine
    \For{$i \gets 1$ \KwTo $n$}{
        $m[i,i] \gets 0$ \tcp{初始化对角线}
    }
    \For{$l \gets 2$ \KwTo $n$}{ \tcp{$l$为链长度}
        \For{$i \gets 1$ \KwTo $n-l+1$}{
            $j \gets i + l - 1$\;
            $m[i,j] \gets \infty$\;
            \For{$k \gets i$ \KwTo $j-1$}{
                $q \gets m[i,k] + m[k+1,j] + p[i-1] \times p[k] \times p[j]$\;
                \If{$q < m[i,j]$}{
                    $m[i,j] \gets q$\;
                    $s[i,j] \gets k$ \tcp{记录最优分割点}
                }
            }
        }
    }
    \caption{MatrixChainOrder}
\end{algorithm}

\begin{algorithm}[t]
    \KwIn{分割矩阵$s[1..n-1,2..n]$, 矩阵编号$i,j$}
    \KwOut{最优括号化方案}
    \BlankLine
    \If{$i == j$}{
        输出 "$A_i$"\;
    }
    \Else{
        输出 "(" \;
        PrintOptimalParens($s$, $i$, $s[i,j]$)\;
        PrintOptimalParens($s$, $s[i,j]+1$, $j$)\;
        输出 ")" \;
    }
    \caption{PrintOptimalParens}
\end{algorithm}

算法的时间复杂度是$\Theta(n^3)$，空间复杂度是$\Theta(n^2)$

\subsection{多源最短路径问题}
\begin{algorithm}[t]
    \KwIn{$n\times n$维矩阵$l[1..n,1..n]$，对于有向图$G = (\{1,2,\ldots,n\},E)$中的边$(i,j)$的长度是$l[i,j]$}
    \KwOut{矩阵$D$使得$D[i,j]=i$到$j$的距离}
    \BlankLine
    $D \leftarrow l$\; {将输入矩阵$l$复制到$D$}
    \For{$k \leftarrow 1$ \KwTo $n$}{
        \For{$i \leftarrow 1$ \KwTo $n$}{
            \For{$j \leftarrow 1$ \KwTo $n$}{
                $D[i,j] = \min\{D[i,j], D[i,k] + D[k,j]\}$\;
            }
        }
    }
    \caption{FLOYD算法}
\end{algorithm}

\subsection{0/1 背包问题}
\begin{align*}
    &\max \sum_{i=1}^nv_ix_i\\
    &s.t.\begin{cases}
        $$
            \sum_{i=1}^nw_ix_i \leq C\\
            x_i \in \{0,1\}\quad i=1,\cdots,n
        $$
    \end{cases}
\end{align*}

设$(y_1,y_2,\cdots,y_n)$是上式的一个最优解，下证$(y_2,y_3,\cdots,y_n)$是下面子问题的最优解：
\begin{align*}
    &\max \sum_{i=2}^nv_ix_i\\
    &s.t.\begin{cases}
        $$
            \sum_{i=2}^nw_ix_i \leq C - w_1y_1\\
            x_i \in \{0,1\}\quad i=2,\cdots,n
        $$
    \end{cases}
\end{align*}

考虑反证，假设$(z_2,z_3,\cdots,z_n)$是上述子问题的最优解，则：
\begin{align*}
    \sum_{i=2}^nv_iz_i\ge\sum_{i=2}^nv_iy_i\\
    w_1y_1+\sum_{i=2}^nw_iz_i\le C\\
\end{align*}

所以有
\[
    v_1y_1+\sum_{i=2}^nw_iz_i\ge \sum_{i=1}{n}v_iy_i
\]

这说明$(y_1,z_2,z_3,\cdots,z_n)$是比$(y_1,y_2,\cdots,y_n)$更优的一个解，矛盾！

背包问题的最优解能够在$\Theta(nC)$的时间内和$\Theta(C)$的空间内得到。

\begin{algorithm}[t]
    \KwIn{物品集合$U=\{u_1,u_2,\ldots,u_n\}$, 体积$s_1,s_2,\ldots,s_n$, 价值$v_1,v_2,\ldots,v_n$, 背包容量$C$}
    \KwOut{$\sum_{u_i\in S}v_i$的最大总价值，且满足$\sum_{u_i\in S}s_i \leq C$，其中$S \subseteq U$}
    \BlankLine
    \For{$i \leftarrow 0$ \KwTo $n$}{
        $V[i,0] \leftarrow 0$\;
    }
    \For{$j \leftarrow 0$ \KwTo $C$}{
        $V[0,j] \leftarrow 0$\;
    }
    \For{$i \leftarrow 1$ \KwTo $n$}{
        \For{$j \leftarrow 1$ \KwTo $C$}{\tcp{如果是滚动数组的话，这里需要是倒序的}
            $V[i,j] \leftarrow V[i-1,j]$\;
            \If{$s_i \leq j$}{
                $V[i,j] \leftarrow \max\{V[i,j], V[i-1,j-s_i]+v_i\}$\;
            }
        }
    }
    \Return $V[n,C]$\;
    \caption{KNAPSACK算法}
\end{algorithm}


\subsection{图像压缩问题}

对于给定的$n$个像素点$p_1, p_2, p_3, \ldots, p_n$，将其分割为连续的$m$个段$S_1, S_2, S_3, \ldots, S_m$。设第$i$个像素段$S_i$中，有$N[i]$个像素（$1 \leq i \leq m$），每个像素用$b[i]$位表示（比如$0$-$15$的数字，要用$4$位表示）。

设$h_i$为分段$S_i$中最大灰度值所对应的二进制数的位数，即
$$
h_{i} = \left\lceil \log \left( \max_{t[i]+1 \leq k \leq t[i]+N[i]} p_{k} + 1 \right) \right\rceil
$$

由于$p_k \leq 255$，则$h_i \leq b[i] \leq 8$，需要用$3$位来表示$b[i]$。我们给定$N[i]$最多是255，这样需要8位来存储。像素段$S_i$需要的总内存空间是$N[i] * b[i] + 11$。

像素序列$\{p_1, p_2, p_3, \ldots, p_n\}$所需存储空间是：

$$
\sum_{i=1}^{m} N[i] \times b[i] + 11m
$$

假设图像压缩问题\textbf{不}满足最优子结构性质，即存在某个最优分段方案，其中至少一个子段的分割不是局部最优的。

考虑像素序列$p_1$到$p_n$的最优分割方案$S_1,\ldots,S_m$，设其中子段$S_k$（$1 \leq k \leq m$）不是该子区间的最优分割。则存在更优的分割$S_k^*$使得：
\begin{equation*}
N[k] \times b[k] + 11 > N^*[k] \times b^*[k] + 11
\end{equation*}

将原方案中的$S_k$替换为$S_k^*$，得到新方案的总存储空间：
\begin{align*}
\text{Total}^* &= \sum_{i\neq k} (N[i] \times b[i] + 11) + (N^*[k] \times b^*[k] + 11) \\
&< \sum_{i=1}^{m} (N[i] \times b[i] + 11) = \text{Total}
\end{align*}
这与原方案的最优性矛盾。因此，最优分割方案中的每个子段必须自身是最优分割。

定义函数 $bm(i,j)$ 表示从像素$i$到$j$中像素灰度值$p_k$最大值所需的最少二进制位数：
\begin{equation*}
bm(i,j) = \left\lceil\log \left(\max_{i \leq k \leq j}\{p_k\} + 1\right)\right\rceil \quad (1 \leq i \leq n)
\end{equation*}
根据动态规划原理，存储长度满足：
\begin{equation}
S[i] = \min_{1 \leq k \leq \min\{i, 256\}} \Big\{ S[i-k] + k \times bm(i-k+1,i)\Big\} + 11
\end{equation}
\begin{lstlisting}
//s[i]记录1-i像素所需的最小空间，I[i]记录1-i像素的最优划分的最后一个段的大小，
//b[i]由p[i]得出，它记录p[i]需要几个比特来存储。
void compress2(int* p, int n, int* s, int* I, int* b)
{
	int header = 11;//一个段需要3位来存储每个像素点需要几个比特来存储，8位来存储这个像素段有多少个像素点。
	int Imax = 256;//N[i]的最大值
	s[0] = 0;
        //这个循环是为了计算s[i]数组，用I[i]数组保存在哪里分段的信息，其它数组、变量都是辅助。
	for (int i = 1; i <= n; i++) {
		//注意s[i]初始化为INT_MAX，而不是初始化为0，因为s[i]要寻找最小值
		b[i] = length(p[i]);
		int bmax = 0;//我们要算出在一个像素段中，最大的那个像素点需要用几个比特来存储，然后这个像素段中所有点都用那么多比特来存储。
		s[i]=INT_MAX;//s[i]求的是最小值，所以初始化为大整数。
		I[i];
		for (int k = 1; k <= i && k <= Imax; k++) {
			if (bmax < b[i - k + 1]) bmax = b[i - k + 1];//如果像素段出现更大像素点，需要更大存储空间，刷新bmax
			if (s[i] > s[i - k] + k * bmax) {//出现更小的s[i]则刷新
				s[i] = s[i - k] + k * bmax;
				I[i] = k;记录分割位置：第1-i像素点的最后一个段有多少数字。
			}
		}
		s[i] += header;
	}
}
\end{lstlisting}

时间复杂度$O(n)$

\subsection{最大子段和}
\begin{lstlisting}
cin>>n;
for(int i=0;i<n;i++){
    cin>>a;
    f=max(f+a,a);
    ans=max(f,ans);
}
\end{lstlisting}
\subsection{多段图问题}
\begin{itemize}
    \item $P(i,j)$: 从$V_i$中结点$j$到汇点$t$的最短路径
    \item $\operatorname{COST}(i,j)$: $P(i,j)$的路径成本
    \item $D(i,j)$: $V_i$中结点$j$的最优决策（选择的后继节点）
\end{itemize}

$$
\operatorname{COST}(i, j)=
\begin{cases} 
0 & \text{当 } i=k \text{（终止条件）} \\ 
\min\limits_{\substack{l \in V_{i+1} \\ \langle j, l \rangle \in E}} 
\big\{ c(j, l) + \operatorname{COST}(i+1, l) \big\} & \text{当 } i<k \text{（递推规则）}
\end{cases}
$$

\begin{itemize}
    \item $BP(i,j)$: 从起点到阶段$V_i$节点$j$的最短路径
    \item $\operatorname{BCOST}(i,j)$: 该路径的累计成本
    \item $D(i,j)$: 节点$j$选择的最优前驱节点
\end{itemize}

$$
\operatorname{BCOST}(i, j)=
\begin{cases} 
0 & \text{基础情况（起点本身）} \\ 
\min\limits_{\substack{l \in V_{i-1} \\ \text{边}(l,j) \in E}} 
\Big[ \operatorname{BCOST}(i-1, l) + c(l, j) \Big] & \text{递推情况}
\end{cases}
$$

\begin{itemize}
    \item 将$n$个资源分配给$r$个项目
    \item 分配$j$个资源给项目$i$时，获得净利$N(i,j)$
    \item 目标：找到使总净利最大的分配方案
\end{itemize}

\begin{enumerate}
    \item \textbf{阶段划分}：
    \begin{itemize}
        \item 构造$r+1$段有向图（$r$个项目对应$r$个决策阶段+1个终止阶段）
        \item 段$i$（$1 \leq i \leq r$）表示项目$i$的分配决策
    \end{itemize}
    
    \item \textbf{节点设计}：
    \begin{align*}
        &\bullet\ \text{源节点 } s = V(1,0) \\
        &\bullet\ \text{中间节点 } V(i,j) = 
        \begin{cases}
            \text{共分配$j$个资源给项目$1,\dots,i-1$}, & 2 \leq i \leq r \\
            j \in \{0,1,\dots,n\}
        \end{cases} \\
        &\bullet\ \text{汇节点 } t = V(r+1,n)
    \end{align*}
    
    \item \textbf{边构造规则}：
    \begin{itemize}
        \item 边形式：$\langle V(i,j), V(i+1,l) \rangle$，满足$j \leq l$且$1 \leq i \leq r$
        \item 边权重：
        $$
        w(\langle V(i,j), V(i+1,l) \rangle) = 
        \begin{cases}
            N(i,l-j), & 1 \leq i < r \\
            \max \{ N(r,p) \mid 0 \leq p \leq n-j \}, & i = r
        \end{cases}
        $$
    \end{itemize}
    
    \item \textbf{求解方法}：
    \begin{itemize}
        \item 最优解对应$s$到$t$的\textbf{最大成本路径}
        \item 通过将边权重取负（$-N(i,k)$）转换为多段图最短路径问题
        \item 动态规划递推式：
        $$
        \operatorname{COST}(i,j) = \max_{\substack{l \geq j \\ \langle V(i,j), V(i+1,l) \rangle \in E}} \big\{ w(\langle V(i,j), V(i+1,l) \rangle) + \operatorname{COST}(i+1,l) \big\}
        $$
    \end{itemize}
\end{enumerate}

\subsection{最优二分检索树问题}
通过动态规划选择树结构，使得：
\begin{itemize}
    \item 成功检索（命中关键字）和
    \item 不成功检索（未命中）的\textbf{加权平均比较次数最少}
\end{itemize}

\begin{equation*}
    \min \left( \sum_{i=1}^{n} p_i \cdot \left( \operatorname{depth}(a_i)+1 \right) + \sum_{j=0}^{n} q_j \cdot \operatorname{depth}(E_j) \right)
\end{equation*}

\begin{tabular}{ll}
    $p_i$ & : 关键字$a_i$的命中概率 \\
    $q_j$ & : 区间$E_j$（失败检索区间）的未命中概率 \\
    $\operatorname{depth}(\cdot)$ & : 结点在树中的深度 \\
\end{tabular}

\begin{equation*}
    C(i,j) = \min_{i \leq k \leq j} \left\{ C(i,k-1) + C(k+1,j) + W(i,j) \right\}
\end{equation*}

\begin{algorithm}[t]
    \KwIn{成功概率数组 $p[1..n]$, 失败概率数组 $q[0..n]$, 关键字数量 $n$}
    \KwOut{期望代价矩阵 $e[1..n+1][0..n]$, 根节点矩阵 $root[1..n][1..n]$}
    \BlankLine
    \For{$i \gets 1$ \KwTo $n+1$}{
        $e[i][i-1] \gets q[i-1]$\;
        $w[i][i-1] \gets q[i-1]$\;
    }
    \For{$\ell \gets 1$ \KwTo $n$}{
        \For{$i \gets 1$ \KwTo $n-\ell+1$}{
            $j \gets i + \ell - 1$\;
            $e[i][j] \gets \infty$\;
            $w[i][j] \gets w[i][j-1] + p[j] + q[j]$\;
            \For{$r \gets i$ \KwTo $j$}{
                $t \gets e[i][r-1] + e[r+1][j] + w[i][j]$\;
                \If{$t < e[i][j]$}{
                    $e[i][j] \gets t$\;
                    $root[i][j] \gets r$\;
                }
            }
        }
    }
    \Return $(e, root)$\;
    \caption{Optimal-BST}
\end{algorithm}

\subsection{投资问题}
\begin{algorithm}[t]
    \KwIn{效益函数 $f[1..n][0..m]$, 总资金 $m$, 项目数 $n$}
    \KwOut{最大效益值 $F[n][m]$, 最优分配方案 $alloc[1..n]$}
    \BlankLine
    \tcp{初始化阶段：处理单个项目情况}
    \For{$x \gets 0$ \KwTo $m$}{
        $F[1][x] \gets f[1][x]$\;
        $tempAlloc[1][x] \gets x$\; \tcp{记录初始分配方案}
    }
    \BlankLine
    \tcp{动态规划递推阶段}
    \For{$k \gets 2$ \KwTo $n$}{
        \For{$x \gets 0$ \KwTo $m$}{
            $F[k][x] \gets 0$\;
            \For{$x_k \gets 0$ \KwTo $x$}{
                $current \gets f[k][x_k] + F[k-1][x - x_k]$\;
                \If{$current > F[k][x]$}{
                    $F[k][x] \gets current$\;
                    $tempAlloc[k][x] \gets x_k$\; \tcp{保存当前最优分配}
                }
            }
        }
    }
    \BlankLine
    \tcp{回溯最优分配方案}
    $remaining \gets m$\;
    \For{$k \gets n$ \textbf{downto} $1$}{
        $alloc[k] \gets tempAlloc[k][remaining]$\;
        $remaining \gets remaining - alloc[k]$\;
    }
    \BlankLine
    \Return $(F[n][m], alloc)$\;
    \caption{Investment-Allocation-DP}
\end{algorithm}

给定$m$元资金和$n$个投资项目，每个项目的效益函数为$f_i(x)$，表示投入$x$元到第$i$个项目的收益。目标是通过资金分配最大化总效益：

\begin{equation*}
    \max \sum_{i=1}^n f_i(x_i) \quad \text{s.t.} \quad \sum_{i=1}^n x_i = m,\ x_i \in \mathbb{N}
\end{equation*}

定义状态函数：
\begin{equation*}
    F_k(x) = \text{将$x$元分配给前$k$个项目的最大效益}
\end{equation*}

递推关系与边界条件：
\begin{align*}
    F_k(x) &= \max_{0 \leq x_k \leq x} \left\{ f_k(x_k) + F_{k-1}(x - x_k) \right\}, \quad k > 1 \\
    F_1(x) &= f_1(x)
\end{align*}

\subsection{本章作业}
\subsubsection{改进floyd算法}
请修改求解所有点对的Floyd算法，使它可以检测到负权圈（一个负权圈是它的总长度为负的圈）。

\begin{itemize}
    \item \textbf{负权圈问题}：当图中存在总权重为负的环路时，Floyd算法会陷入无限更新状态。例如，若路径$i \rightarrow j$包含负权圈，每次绕行该圈都会使总路径权重减小，导致算法无法收敛。
    
    数学表现为：
    $$
    \exists \text{环 } C: \sum_{e \in C} w(e) < 0 \Rightarrow \lim_{k \to \infty} dist^{(k)}[i][j] = -\infty
    $$
    
    \item \textbf{最优子结构破坏}：动态规划要求子问题的最优解能构成全局最优解。负权圈使得通过无限绕行可不断降低路径权重，导致：
    \begin{align*}
        \text{局部最优} &\nrightarrow \text{全局最优} \\
        dist[i][j] &= \min(dist[i][j], dist[i][k] + dist[k][j]) \quad \text{（不收敛）}
    \end{align*}
\end{itemize}

\textbf{改进思路}：在算法完成后检查距离矩阵对角线元素，满足：
$$
\exists i \in V: dist[i][i] < 0 \Rightarrow \text{存在包含顶点$i$的负权圈}
$$

\subsubsection{有向图的所有点对最短路径问题}
设G是一个有向图，其中每一条边有一个非负的长度，如果两点之间无边，令其长度为无穷。求图中每个顶点到其他顶点的最短路径。请证明该问题满足最优子结构性质。

假设最短路径问题不满足最优子结构性质，则存在以下情形：

\begin{enumerate}
    \item \textbf{假设}：
    设$p = \langle v_{1},v_{2},\ldots,v_{k} \rangle$是从$v_{1}$到$v_{k}$的最短路径，但存在子路径$p_{ij} = \langle v_{i},v_{i+1},\ldots,v_{j} \rangle$不是$v_{i}$到$v_{j}$的最短路径。
    
    \item \textbf{构造矛盾}：
    \begin{itemize}
        \item 存在更优路径$p'_{ij}$满足$w(p'_{ij}) < w(p_{ij})$
        \item 构造新路径$p' = \langle v_{1},\ldots,v_{i-1} \rangle \circ p'_{ij} \circ \langle v_{j+1},\ldots,v_{k} \rangle$
        \item 计算总权重：
        $$
        w(p') = w(p) - w(p_{ij}) + w(p'_{ij}) < w(p)
        $$
    \end{itemize}
    
    \item \textbf{矛盾结论}：
    这与$p$是最短路径的初始假设矛盾，故原命题得证。
\end{enumerate}


\subsubsection{k-shattering 问题}
\paragraph{问题定义}
给定边带权值的完全二叉树$T=(V,E,w)$，其中$w: E \rightarrow \mathbb{R}^+$。一个\emph{k-shattering集合}是边集$S \subseteq E$满足：
\begin{itemize}
    \item 删除$S$后，$T$被划分为$m$棵子树$\{T_1,...,T_m\}$
    \item 每棵子树$T_i$的节点数$|V(T_i)| \leq k$
\end{itemize}
求使$\sum_{e\in S}w(e)$最小的k-shattering集合。


\begin{algorithm}[H]
\SetKwFunction{kshatter}{kShatteringDP}
\SetKwProg{Fn}{Function}{}{}
\Fn{\kshatter{T, w, k}}{
    $n \gets |V(T)|$\;
    \texttt{dp}[1..n][1..k] $\gets \infty$\;
    
    \For{$i \gets n$ \textbf{downto} $\lfloor n/2 \rfloor +1$}{
        \For{$m \gets 1$ \textbf{to} $k$}{
            $\texttt{dp}[i][m] \gets 0$ \tcp*{叶子节点初始化}
        }
    }
    
    \For{$i \gets \lfloor n/2 \rfloor$ \textbf{downto} $1$}{
        \For{$j \gets 1$ \textbf{to} $k$}{
            \uIf{$j = 1$}{
                $\texttt{min\_left} \gets \min\limits_{1\leq a \leq k}\texttt{dp}[2i][a]$\;
                $\texttt{min\_right} \gets \min\limits_{1\leq b \leq k}\texttt{dp}[2i+1][b]$\;
                $\texttt{dp}[i][1] \gets w(i,2i) + w(i,2i+1) + \texttt{min\_left} + \texttt{min\_right}$\;
            }
            \Else{
                $\texttt{option1} \gets w(i,2i) + \min\limits_{1\leq a \leq k}\texttt{dp}[2i][a] + \texttt{dp}[2i+1][j-1]$\;
                $\texttt{option2} \gets w(i,2i+1) + \min\limits_{1\leq a \leq k}\texttt{dp}[2i+1][a] + \texttt{dp}[2i][j-1]$\;
                $\texttt{option3} \gets \min\limits_{1\leq a \leq j-2}(\texttt{dp}[2i][a] + \texttt{dp}[2i+1][j-a-1])$\;
                $\texttt{dp}[i][j] \gets \min(\texttt{option1}, \texttt{option2}, \texttt{option3})$\;
            }
        }
    }
    \Return $\min\limits_{1\leq j \leq k}\texttt{dp}[1][j]$\;
}
\caption{最小代价k-shattering算法}
\end{algorithm}

\begin{itemize}
    \item \textbf{状态定义}：$\texttt{dp}[i][m]$表示以$i$为根的子树保留$m$个节点的最小代价
    \item \textbf{边界条件}：
    $$
    \texttt{dp}[i][m] = 0 \quad \forall i > \lfloor n/2 \rfloor
    $$
    \item \textbf{状态转移}：
    $$
    \texttt{dp}[i][j] = \begin{cases}
    w_{i,2i}+w_{i,2i+1}+\min\limits_{a}\texttt{dp}[2i][a]+\min\limits_{b}\texttt{dp}[2i+1][b]\text{(只有根)} & j=1 \\
    \min\begin{cases}
    w_{i,2i}+\min\texttt{dp}[2i][a]+\texttt{dp}[2i+1][j-1] \text{（断左）}\\
    w_{i,2i+1}+\min\texttt{dp}[2i+1][a]+\texttt{dp}[2i][j-1]\text{（断右）} \\
    \min\limits_{a}(\texttt{dp}[2i][a]+\texttt{dp}[2i+1][j-a-1])\text{（全留）}
    \end{cases} & j>1
    \end{cases}
    $$
    \item \textbf{时间复杂度}：$O(nk^2)$
\end{itemize}

最小代价：$\min(16,9)=9$，对应删除边集$\{\langle b,d\rangle, \langle b,e\rangle, \langle a,c\rangle\}$。

\end{document}